
        \documentclass[fleqn]{article}      
        \usepackage{amsmath}
        \usepackage{fontspec}
        \usepackage{kotex} % 한국어 지원  

        \begin{document}      
        ```latex
\documentclass{article}
\usepackage{amsmath}
\usepackage{geometry}
\geometry{a4paper, margin=1in}
\begin{document}

\section*{고등수학 (상) - 방정식과 부등식 - 연립이차방정식}

\subsection*{문제 1 (쉬움, 연립이차방정식)}
다음 연립이차방정식을 푸시오.
\[
\begin{cases}
x^2 + y^2 = 25 \\
x + y = 7
\end{cases}
\]
\begin{enumerate}
    \item (3,4) 
    \item (4,3) 
    \item (5,2) 
    \item (6,1) 
    \item (7,0)
\end{enumerate}

\subsection*{문제 2 (쉬움, 연립이차방정식)}
다음 연립이차방정식을 푸시오.
\[
\begin{cases}
x^2 + y^2 = 10 \\
x - y = 2
\end{cases}
\]
\begin{enumerate}
    \item (3,1) 
    \item (2,2) 
    \item (1,3) 
    \item (4,0) 
    \item (0,4)
\end{enumerate}

\subsection*{문제 3 (쉬움, 연립이차방정식)}
다음 연립이차방정식을 푸시오.
\[
\begin{cases}
x^2 + y^2 = 16 \\
x + y = 4
\end{cases}
\]
\begin{enumerate}
    \item (2,2) 
    \item (3,1) 
    \item (4,0) 
    \item (0,4) 
    \item (1,3)
\end{enumerate}

\subsection*{문제 4 (쉬움, 연립이차방정식)}
다음 연립이차방정식을 푸시오.
\[
\begin{cases}
x^2 + y^2 = 13 \\
x - y = 1
\end{cases}
\]
\begin{enumerate}
    \item (3,2) 
    \item (2,3) 
    \item (4,1) 
    \item (1,4) 
    \item (0,3)
\end{enumerate}

\subsection*{문제 5 (쉬움, 연립이차방정식)}
다음 연립이차방정식을 푸시오.
\[
\begin{cases}
x^2 + y^2 = 5 \\
x + y = 3
\end{cases}
\]
\begin{enumerate}
    \item (2,1) 
    \item (1,2) 
    \item (3,0) 
    \item (0,3) 
    \item (4,-1)
\end{enumerate}

\subsection*{문제 6 (중간 난이도, 연립이차방정식)}
다음 연립이차방정식을 푸시오.
\[
\begin{cases}
x^2 + y^2 = 20 \\
xy = 12
\end{cases}
\]
\begin{enumerate}
    \item (4,2) 
    \item (3,3) 
    \item (5,1) 
    \item (2,4) 
    \item (0,5)
\end{enumerate}

\subsection*{문제 7 (중간 난이도, 연립이차방정식)}
다음 연립이차방정식을 푸시오.
\[
\begin{cases}
x^2 - y^2 = 7 \\
x + y = 5
\end{cases}
\]
\begin{enumerate}
    \item (6,1) 
    \item (4,3) 
    \item (3,2) 
    \item (5,0) 
    \item (2,3)
\end{enumerate}

\subsection*{문제 8 (중간 난이도, 연립이차방정식)}
다음 연립이차방정식을 푸시오.
\[
\begin{cases}
x^2 + y^2 = 30 \\
x - y = 4
\end{cases}
\]
\begin{enumerate}
    \item (5,1) 
    \item (6,2) 
    \item (7,3) 
    \item (8,0) 
    \item (4,4)
\end{enumerate}

\subsection*{문제 9 (중간 난이도, 연립이차방정식)}
다음 연립이차방정식을 푸시오.
\[
\begin{cases}
x^2 + y^2 = 50 \\
xy = 24
\end{cases}
\]
\begin{enumerate}
    \item (6,4) 
    \item (5,5) 
    \item (8,2) 
    \item (10,0) 
    \item (2,8)
\end{enumerate}

\subsection*{문제 10 (중간 난이도, 연립이차방정식)}
다음 연립이차방정식을 푸시오.
\[
\begin{cases}
x^2 - y^2 = 8 \\
x + y = 6
\end{cases}
\]
\begin{enumerate}
    \item (5,1) 
    \item (4,2) 
    \item (6,0) 
    \item (3,3) 
    \item (2,4)
\end{enumerate}

\subsection*{문제 11 (어려움, 연립이차방정식)}
연립이차방정식
\[
\begin{cases}
x^2 + y^2 = 45 \\
xy = 20
\end{cases}
\]
을 만족하는 \( (x, y) \)의 한 쌍을 구하시오.

\subsection*{문제 12 (어려움, 연립이차방정식)}
연립이차방정식
\[
\begin{cases}
x^2 + y^2 = 34 \\
x - y = 2
\end{cases}
\]
을 만족하는 \( (x, y) \)의 한 쌍을 구하시오.

\subsection*{문제 13 (어려움, 연립이차방정식)}
연립이차방정식
\[
\begin{cases}
x^2 - y^2 = 15 \\
xy = 10
\end{cases}
\]
을 만족하는 \( (x, y) \)의 한 쌍을 구하시오.

\subsection*{문제 14 (어려움, 연립이차방정식)}
연립이차방정식
\[
\begin{cases}
x^2 + y^2 = 20 \\
x + 2y = 10
\end{cases}
\]
을 만족하는 \( (x, y) \)의 한 쌍을 구하시오.

\subsection*{문제 15 (어려움, 연립이차방정식)}
연립이차방정식
\[
\begin{cases}
x^2 + 2y^2 = 24 \\
x + y = 8
\end{cases}
\]
을 만족하는 \( (x, y) \)의 한 쌍을 구하시오.

\section*{정답 및 해설}

\subsection*{문제 1 해설}
주어진 연립이차방정식을 대입법을 사용하여 풀면,  
\[
(x+y)^2 - 2xy = 25
\]
\[
49 - 2xy = 25
\]
\[
xy = 12
\]
따라서 정답은 \( \boxed{(3,4)} \) 또는 \( \boxed{(4,3)} \)이다.

\subsection*{문제 2 해설}
주어진 연립이차방정식을 대입법을 사용하여 풀면,  
\[
x = y + 2 \Rightarrow (y + 2)^2 + y^2 = 10
\]
\[
2y^2 + 4y - 6 = 0 \Rightarrow y^2 + 2y - 3 = 0
\]
따라서 정답은 \( \boxed{(3,1)} \) 또는 \( \boxed{(1,3)} \)이다.

\subsection*{문제 3 해설}
주어진 연립이차방정식을 대입법을 사용하여 풀면,  
\[
x + y = 4 \Rightarrow y = 4 - x
\]
\[
x^2 + (4 - x)^2 = 16
\]
\[
x^2 + 16 - 8x + x^2 = 16 \Rightarrow 2x^2 - 8x = 0
\]
따라서 정답은 \( \boxed{(2,2)} \)이다.

\subsection*{문제 4 해설}
주어진 연립이차방정식을 대입법을 사용하여 풀면,  
\[
y = x - 1
\]
\[
x^2 + (x - 1)^2 = 13
\]
\[
2x^2 - 2x - 12 = 0 \Rightarrow x^2 - x - 6 = 0
\]
따라서 정답은 \( \boxed{(3,2)} \) 또는 \( \boxed{(2,3)} \)이다.

\subsection*{문제 5 해설}
주어진 연립이차방정식을 대입법을 사용하여 풀면,  
\[
y = 3 - x
\]
\[
x^2 + (3 - x)^2 = 5
\]
\[
2x^2 - 6x + 4 = 0 \Rightarrow x^2 - 3x + 2 = 0
\]
따라서 정답은 \( \boxed{(2,1)} \) 또는 \( \boxed{(1,2)} \)이다.

\subsection*{문제 11 해설}
주어진 연립이차방정식을 대입법을 사용하여 풀면,
\[
x^2 + y^2 = 45, \quad xy = 20
\]
\[
x^2 + y^2 = (x+y)^2 - 2xy \Rightarrow (x+y)^2 = 45 + 40 = 85
\]
따라서 \( x+y = \sqrt{85} \) 또는 \( x+y = -\sqrt{85} \)이다.

\subsection*{문제 12 해설}
주어진 연립이차방정식을 대입법을 사용하여 풀면,
\[
y = x - 2
\]
\[
x^2 + (x - 2)^2 = 34
\]
\[
2x^2 - 4x - 30 = 0 \Rightarrow x^2 - 2x - 15 = 0
\]
따라서 정답은 \( \boxed{(5,3)} \) 또는 \( \boxed{(3,5)} \)이다.

\subsection*{문제 13 해설}
주어진 연립이차방정식을 대입법을 사용하여 풀면,
\[
x^2 - y^2 = (x+y)(x-y) = 15, \quad xy = 10
\]
\[
x+y = m, \quad x-y = n \Rightarrow mn = 15
\]
따라서 정답은 \( \boxed{(5,2)} \) 또는 \( \boxed{(2,5)} \)이다.

\subsection*{문제 14 해설}
주어진 연립이차방정식을 대입법을 사용하여 풀면,
\[
y = 10 - x
\]
\[
x^2 + 2(10 - x)^2 = 20
\]
\[
x^2 + 2(100 - 20x + x^2) = 20
\]
따라서 정답은 \( \boxed{(4,6)} \) 또는 \( \boxed{(6,4)} \)이다.

\subsection*{문제 15 해설}
주어진 연립이차방정식을 대입법을 사용하여 풀면,
\[
y = 8 - x
\]
\[
x^2 + 2(8 - x)^2 = 24
\]
\[
x^2 + 2(64 - 16x + x^2) = 24
\]
따라서 정답은 \( \boxed{(8,0)} \) 또는 \( \boxed{(0,8)} \)이다.

\end{document}
```      
        \end{document} 
    